\section{Joint majorants for Gaussian prepared factors}
\label{supp:joint-majorants}

The factors produced by the local change of variables share one Gaussian
law.  Estimates for separate expectations do not suffice.  The following
estimate is the product bound used before the connected expansion.  Its constants do not depend on the dimension of that Gaussian.

\begin{lemma}[A joint Gaussian product estimate]
\label{supp:gaussian-product}
Let $Z$ be a centered Gaussian vector with covariance $C$, where
$0\le C\le c_*I$ and $c_*>0$.  Let $O_i$ satisfy $|O_i|\le a_i$.
For a finite collection of ratio factors put
\[
 R_\alpha=e^{Q_\alpha}-1,\qquad
 |Q_\alpha(Z)|\le e_\alpha
       \left(1+\sum_{x\in I_\alpha}|Z_x|^2\right),
 \qquad 0<e_\alpha\le1.
\]
Suppose that
\begin{align}
 \sup_x\sum_{\alpha:x\in I_\alpha}\sqrt{e_\alpha}
       &\le(64c_*)^{-1},\label{supp:endpoint-load}\\
 (2+4c_*|I_\alpha|)\sqrt{e_\alpha}
       &\le\tfrac14\log(e_\alpha^{-1})
       \quad\hbox{for every }\alpha.\label{supp:ratio-size}
\end{align}
Let $F_1,\ldots,F_n$ be bounded factors whose product has absolute
value at most $\mathbf1_E$, where
$\mathbb P(E)\le e^{-\tau n}$.  Finally suppose $\|S\|_{L^4}\le B$.
Then
\[
 \mathbb E\left|S\prod_iO_i\prod_\alpha R_\alpha
                         \prod_{j=1}^nF_j\right|
 \le B\prod_i a_i\prod_\alpha e_\alpha^{1/4}
                                      e^{-\tau n/2}.
\]
The convention for $n=0$ is $E$ equal to the full probability space.
The conclusion also holds for factors $R_\alpha$ satisfying directly
$|R_\alpha|\le\sqrt{e_\alpha}\exp\{2\sqrt{e_\alpha}
(1+\sum_{x\in I_\alpha}|Z_x|^2)\}$.
Here $x$ indexes scalar Gaussian coordinates; vector coordinates may be
replaced by their scalar components.
No independence between any of the displayed factors is required.
\end{lemma}

\begin{proof}
For $U\ge1$ and $0<e\le1$, the elementary exponential estimate gives
\[
 |e^Q-1|\le eUe^{eU}\le\sqrt e\,e^{2\sqrt e\,U}
 \quad\text{if }|Q|\le eU.
\]
Let $D$ be the diagonal matrix whose $x$th entry is
$\sum_{\alpha:x\in I_\alpha}\sqrt{e_\alpha}$.
The Gaussian determinant identity, valid also for a singular covariance
by approximation, yields
\[
 \mathbb E e^{8Z^TDZ}
 =\det(I-16C^{1/2}DC^{1/2})^{-1/2}
 \le e^{16c_*\operatorname{tr}D}.
\]
Indeed, \eqref{supp:endpoint-load} makes the eigenvalues of
$16C^{1/2}DC^{1/2}$ at most $1/4$, and
$-\log(1-u)\le2u$ for $0\le u\le1/2$.
Consequently
\[
 \left\|\prod_\alpha R_\alpha\right\|_{L^4}
 \le\prod_\alpha\left[\sqrt{e_\alpha}
       e^{(2+4c_*|I_\alpha|)\sqrt{e_\alpha}}\right]
 \le\prod_\alpha e_\alpha^{1/4}.
\]
H\"older's inequality with exponents $4,4,2$, applied respectively to
$S$, the ratio product, and $\mathbf1_E$, completes the proof.
\end{proof}

Here is the precise reserve bookkeeping used when taking a fixed number
of derivatives.  It separates the Gaussian estimate from the geometric
verification of its hypotheses.

\begin{corollary}[Fixed-order derivatives and an exceptional reserve]
\label{supp:product-reserve}
Let the factors in Lemma~\ref{supp:gaussian-product} carry integer loads
$s\ge1$, and let $s_{\rm tot}$ be their total load.  Fix an integer $k$.
Suppose that each differentiated product, through order $k$, is a sum of
at most $C_k(1+s_{\rm tot})^{d_k}$ expressions to which that lemma
applies.  Each expression retains the original factor indices and
envelopes $a_i,e_\alpha,I_\alpha$, the same failure count $n$, and the
same core loads; differentiated factors may use the direct ratio
envelope stated in that lemma.  Suppose, for every such expression,
\[
 B\le C_k g^{-d_k}(1+s_{\rm tot})^{d_k},\qquad 0<g<1.
\]
Suppose each expression requiring the factor $g^{-d_k}$ has at least
one of the following reserves:
\begin{enumerate}
\item a deterministic core factor bounded by
$e^{-\kappa p^2s_c}$, with $\kappa p^2/2\ge d_k\log(g^{-1})$;
\item at least one failure test, with
$\tau/4\ge d_k\log(g^{-1})$;
\item at least one ratio factor, with
$\log(e_\alpha^{-1})/8\ge d_k\log(g^{-1})$;
\item a deterministic exceptional factor with envelope $a_*\le e_*$,
where $0<e_*<1$ and
$\log(e_*^{-1})/2\ge d_k\log(g^{-1})$.
\end{enumerate}
Expressions without any such factor are assumed to obey the same
bound on $B$ without $g^{-d_k}$.  Then the derivative has a joint
product majorant with deterministic ordinary weights $a_i$, core
weights $e^{-\kappa p^2s_c/2}$, failure weights $e^{-\tau/4}$, and
ratio weights $e_\alpha^{1/8}$, and deterministic exceptional weights
$e_*^{1/2}$ replacing their original envelopes, times a fixed polynomial in
$1+s_{\rm tot}$.  For a single marked parameter or input discrepancy,
factor out its size and use an envelope for the remaining marked factor
in place of its original envelope.  No bound relative to the original
factor is required; that factor may vanish while its discrepancy does
not.  If the hypotheses above, including the reserves, hold with this
replacement, and all other envelopes hold uniformly along the comparison
or interpolation, the same conclusion retains the discrepancy size.
\end{corollary}

\begin{proof}
Apply Lemma~\ref{supp:gaussian-product} to each expression.  Pay the
single global factor $g^{-d_k}$ from the first available reserve in
the fixed order core, failure, ratio, deterministic exceptional.
The four displayed inequalities
justify precisely this payment.  Weakening the unused reserves to the
same final exponents only enlarges the bound.  Sum the polynomially
many derivative placements.  For a marked term apply the same argument
with its replacement envelope.  The discrepancy size remains a common
multiplicative factor and is not used for this payment.
\end{proof}

A fixed polynomial of the total load is harmless for the connected
expansion: for every $\varepsilon>0$ and fixed $d$,
$(1+s)^d\le C_{d,\varepsilon}e^{\varepsilon s}$.  It is therefore paid
from an explicitly reserved part of the support exponent.

To apply these lemmas to the block integration, the geometric
construction supplies the deterministic ordinary and core envelopes,
the simultaneous-failure estimate
$\mathbb P(E)\le\exp\{-c_Lp^2n/M^D\}$, and the endpoint envelopes of
the determinant and Gaussian-ratio chains.  Those chains have
exponential accuracy in $M$ with polynomial endpoint counts; their
summed square roots satisfy \eqref{supp:endpoint-load}, and
\eqref{supp:ratio-size} follows after increasing the terminal scale.
Their eighth-root reserve pays the fixed-order score losses.
The scale choices ensure
\[
 \frac{c_Lp^2}{M^D}\gg p^{3/10}+\log(g^{-1}),
\]
so the failure and core reserves also satisfy
Corollary~\ref{supp:product-reserve} and retain the required
$e^{-p^{3/10}}$ suppression.

For a component containing measurable physical-spin factors, use one
Gaussian marginal on the union of its read coordinates and one
transport vector field that fixes all of those physical-spin arguments.
In addition to the covariance bounds, the construction supplies bounds
through the required fixed derivative order on the local mean, the
transport vector field and its divergence, and, for a varying covariance,
its parameter derivatives.  These bounds follow from the smooth inverse
chart on a fixed compact subset of its injectivity ball and the
corresponding differentiated window-kernel bounds.  They make each
fixed-order transport score bounded in $L^4$ by a polynomial of the load
and a fixed power of $g^{-1}$.  Uniform ellipticity alone would not imply
this score bound.  A product-rule derivative does not remove
the closed support conditions of a retained sector profile.  Thus the
geometric reserve on those supports, proved in
Appendix~\ref{app:rg-geometry}, is exactly the
hypothesis of Corollary~\ref{supp:product-reserve}; no factor-by-factor
independence argument is involved.


\subsection{Verification for the eight classes of prepared factors}
\label{app:joint-application}

We apply Lemma~\ref{supp:gaussian-product} on the union of the read
coordinates of a finite selection in one fixed primary pattern.
The scalar-coordinate covariance bound follows from
\eqref{rg:window-bounds}; it is uniform in the pattern, its exterior,
and the admitted period. Empty exterior means a zero-dimensional
Gaussian and causes no exception. The complete read sets and attached
profile rules are those of Appendix~\ref{app:rg-geometry}.

For the eight classes listed in Section~\ref{rg:prepared}, the following
are the required envelopes and reserves.
\begin{enumerate}
\item The unassigned nonlinear rows have deterministic envelopes
$C_Lg\operatorname{poly}(M,p)$. Their normalized smooth derivatives
retain that order: a cutoff derivative contributes $t^{-1}$ together
with a direction of size $t$ times its prescribed distance weight.
\item The stationary squares have order $g^2$; the
$\mathcal N,\mathcal Z$ remainders have exponential accuracy in $M$.
Their fixed-order derivatives have the same envelopes by
\eqref{rg:prediction} and the truncated free-kernel estimates.
\item A remote canonical term retains
\eqref{rg:old-polynomial-bound}, and a remote regular error retains its
actual derivative norm. A nonremote term is a selected hard letter
with the attached profiles prescribed in Appendix~\ref{app:rg-geometry};
it meets a compulsory core by its recorded nonremote test.
\item A kinetic tail retains every path and mask query. If it is
smooth its exponential envelope suffices. If it is treated as a hard
letter, that same exceptional envelope pays its fixed-order transport
loss. The preparation is $1+\chi_{\rm att}(e^V-1)$, so it introduces
no unpriced cutoff failure.
\item The ratio-chain letters use \eqref{rg:ratio-envelope}.
Their scalar endpoint counts are polynomial in their recorded loads and
in $M$. To sum their eighth roots, fix a support exponent $B>0$, a
support hit, and a number $r\ge1$ of residual hops. The finite-range convention of
Section~\ref{app:factor-ownership} gives $O(M)$ range for each
matrix hop and window. There are at most $C_LM^D$ choices at each
hop; choosing which of the $O(r)$ retained windows has its fixed
$O(M)$ support neighborhood containing the hit costs only a further
polynomial. Thus the number $N_r$ of literal
paths is at most
\[
 N_r\le\operatorname{poly}(M,r)(C_LM^D)^{r+O(1)}.
\]
For a resolvent chain, first integrate its coefficient in the resolvent
parameter, using the integrable $(1+u)^{-2}$ bound in
Section~\ref{rg:gaussian}. The envelopes $e_\alpha$ then index discrete
paths. Apply \eqref{rg:chain-price} with support exponent $8B$ and
H\"older's inequality for this finite path set to obtain
\[
 \begin{aligned}
 \sum_{\alpha:\,r\text{ hops}}e^{Bs_\alpha}e_\alpha^{1/8}
 &\le N_r^{7/8}
   \left(\sum_{\alpha:\,r\text{ hops}}
                  e^{8Bs_\alpha}e_\alpha\right)^{1/8}\\
 &\le\operatorname{poly}(M,r)
     \exp\{-c_LMr/8+O_L(r\log M)+CBr\}.
 \end{aligned}
\]
This is summable over $r$ and tends to zero as $H$ increases.
Since the envelopes are less than one, their square-root sum has the
same upper bound.
Increasing $H$ makes this support-hit sum less than $(64c_*)^{-1}$.
The exponential decrease in $M$ times the number of hops dominates the
polynomial endpoint count, proving \eqref{supp:ratio-size} as well.
\item The Haar-Jacobian letter is smooth and has order $g^2$ with
its normalized derivatives, since the Jacobian has zero linear term.
\item The opened edge, observation and chart defaults are the failure
factors. Equation~\eqref{rg:joint-rarity} supplies the simultaneous bound
with $\tau=c_Lp^2/M^D$. It holds on all fixed-order derivative supports,
by the prediction estimates and the retained support properties of
Section~\ref{rg:prepared}.
\item An old covering weight retains its sup envelope. Its nonempty
inventory supplies a marked primary input edge, hence a compulsory core.
Its measurable physical-spin arguments are fixed by the common
transport; the weight is never field-differentiated.
\end{enumerate}

The compulsory core itself is \eqref{app:explicit-core}. Its reserve,
proved there on the retained profile and derivative supports, gives
a deterministic envelope $e^{-\kappa_Lp^2s_c}$ after increasing $H$.
Normalized derivatives of its smooth factors or exponent cost only a
fixed polynomial in $g^{-1},M,p$ and the total number of reads.
The same reserve remains during interpolation of two core exponents
at fixed physical spins, because both endpoints obey it.

For a smooth product without a hard factor, the ordinary derivative
bounds already stated preserve the ordinary orders. In a product with
hard physical-spin dependencies, use the single simultaneous transport
at the end of Appendix~\ref{app:rg-geometry}. That calculation supplies
the derivatives of the mean, covariance score, vector field and
divergence required above. There are at most a polynomial number of
placements at a fixed total derivative order. The $L^4$ norm of each
score is bounded by a fixed polynomial of the number of reads and by
fixed powers of $g^{-1},M,p$. Fixed powers of $M,p$ can be absorbed in a
further fixed power of $g^{-1}$ for sufficiently large $H$. Every use
that requires this loss has a core, failure or exceptional ratio/tail
reserve by the eight cases just listed. A marked coupling, history or
input discrepancy retains its own factor in these same bounds, using
the window difference identity and the interpolation at fixed
physical-spin arguments.
For an input discrepancy, telescope a finite product into terms with
one marked factor $F_{\alpha,2}-F_{\alpha,1}$. Its coefficient,
regular-error or covering-weight difference supplies the replacement
envelope in Corollary~\ref{supp:product-reserve}; the other factors
use envelopes uniform for the two inputs. The marked support retains
the same core or exceptional reserve whenever transport requires it.

Corollary~\ref{supp:product-reserve} therefore proves
\eqref{rg:joint-product}. The support-hit sums of the resulting
majorants are obtained exactly as at the end of the proof of
Lemma~\ref{rg:joint}: projected old records use
\eqref{rg:projection}, canonical paths use their fixed exponential
weight, core patterns are paid by their $p^2$ reserve, and chains use
\eqref{rg:chain-price}. Reserve a further fixed positive amount of
support exponent for the polynomial in total load. The conditions
\eqref{rg:finite-choices} leave this reserve while retaining
$e^{-p^{3/10}}$ for output-seeded components. This is a product estimate
under one Gaussian, not an independence assertion about overlapping
factors. Exact independence is used only for disjoint complete supports,
where Lemma~\ref{app:local-marginals} proves it.
